%-------------symbols --------------------------------------
% 
\newcommand\tr{\mathop\mathrm{tr}\nolimits}
\newcommand\Tr{\mathop\mathrm{Tr}\nolimits}
\newcommand\Det{\mathop\mathrm{Det}\nolimits}
\renewcommand\Im{\mathop\mathrm{Im}\nolimits}
%\renewcommand\Re{\mathop\mathrm{Re}\nolimits}
\renewcommand\span{\mathop\mathrm{span}\nolimits}
\newcommand{\sgn}{\mathop{\mathrm{sgn}}\nolimits}
\newcommand\diag{\mathop\mathrm{diag}\nolimits}
%\newcommand\rank{\mathop\mathrm{rank}\nolimits}
\newcommand\Res{\mathop\mathrm{Res}\limits}
\newcommand\eqdef{\buildrel {\rm def}\over=} 
\newcommand\cd{\mathop\mathrm{\kappa}_2\nolimits}
\newcommand\cdest{\mathop\mathrm{\kappa}_2^\mathrm{est}\nolimits}


%--------------colors---------------------------------------
% 
%\newrgbcolor{palered}{1 0.7 0.7}
%\newrgbcolor{paleblue}{0.7 0.7 1}
%\newrgbcolor{darkcyan}{0 0.6 0.6}
%\newrgbcolor{brown}{0.6 0.6 0}
%\newrgbcolor{orange}{1 0.5 0}
%-------------macros--------------------------------------
%
\newcommand\be{\begin{equation}}
\newcommand\ee{\end{equation}}
\newcommand\bea{\begin{eqnarray}}
\newcommand\eea{\end{eqnarray}}
\newcommand\bi{\begin{itemize}}
\newcommand\ei{\end{itemize}}
\newcommand{\odg}[1]{\mathcal{O}\left( #1 \right)}
\newcommand\odgeps{\mathcal{O}(\varepsilon)}
\newcommand\nn{\nonumber}
\newcommand\vev[1]{\left\langle #1\right\rangle}
% ----- ChASE papers -------------
\newcommand\csi{ChASE\xspace}
\newcommand\rto[2]{|\rho_{#1}|^{#2}}
\newcommand\subit[2]{#1^{(#2)}}
\newcommand\her[1]{\left(#1\right)^H}
\newcommand\inv[1]{\left(#1\right)^{-1}}
\newcommand\pl{{\textrm p}_{m}}
\newcommand\p[1]{{\textrm p}_{m_#1}}
\newcommand\rowmat[2]{\left(
  \begin{tabular}{c c} 
    $#1$ & $#2$ 
  \end{tabular}\right)}
\newcommand\colmat[2]{\left(
  \begin{tabular}{c} 
    $#1$ \\ 
    $#2$ 
  \end{tabular}\right)} 
\newcommand\sqrmat[4]{\left(
  \begin{tabular}{c c} 
    $#1$ & $#2$ \\
    $#3$ & $#4$
  \end{tabular}\right)} 
% ----- Special letters -----
\newcommand\Ham{\hat{H}}
\newcommand\T{\hat{T}}
\newcommand\V{\hat{V}}
\newcommand\U{\hat{U}}
\newcommand\W{\hat{W}}
\newcommand\Rcal{\ensuremath{\mathcal{R}}}
\newcommand\Pcal{\ensuremath{\mathcal{P}}}
\newcommand\Scal{\ensuremath{\mathcal{S}}}
\newcommand\Ibb{\ensuremath{\mathbb{1}}} % requires \usepackage{bbold}
\newcommand\nr{{\it n}({\bf r})}
\newcommand\kv{{\bf k}}
\newcommand\rv{{\bf r}}
\newcommand\blk{{\tt blk}}
\newcommand\polym{{\tt deg}}
%\newcommand\nev{{\tt nev}}
\newcommand\sv{$\vec{s}$}
\newcommand\nev{\textsf{nev}}
\newcommand\nex{\textsf{nex}}
% ----- Number fields --------
\newcommand{\N}{\mathbb{N}}
\newcommand{\Z}{\mathbb{Z}}
\newcommand{\Q}{\mathbb{Q}}
\newcommand{\R}{\mathbb{R}}
\newcommand{\C}{\mathbb{C}}
\newcommand{\PP}{\mathbb{P}}
\newcommand{\acosh}{\textrm{cosh}^{-1}}
\newcommand{\seq}{$\{P^{(\ell)}\}$}
\newcommand\argmin{\mathop\mathrm{argmin}\nolimits}
%\let\brk\relax
\newcommand\distimes{$\discretionary{\hbox{${}\times{}\!$}}{}{}$}
\let\VEV=\vev
\newcommand\coeff[2]{{\textstyle{#1\over #2}}}
\newcommand\half{\coeff{1}{2}}
\newcommand\bra[1]{\left\langle #1\right|}
\newcommand\ket[1]{\left| #1\right\rangle}
\newcommand\braket[2]{\VEV{#1 | #2}}
\newcommand\norm[1]{\left\|#1\right\|}
\newcommand\pnorm[2]{\left\|#1\right\|_{#2}}
\newcommand\eq[1]{Eq.~(\ref{#1})}
\newcommand\eqalign[1]{%
	\vcenter{%
		\normalbaselines \advance\baselineskip 5pt
		\advance\lineskip 5pt \tabskip=0pt
		\halign{%
			&\hfil $\displaystyle{##{}}$&
			$\displaystyle{{}##}$\hfil\cr
			#1\crcr
			}%
		}%
	}
\newcommand\eqrange[2]{Eqs.~(\ref{#1}--\reftail{#2})}
\newcommand\dotline{\par\hbox to \hsize{\dotfill}\par}
\newcommand{\hcancel}[1]{%
    \tikz[baseline=(tocancel.base)]{
        \node[inner sep=0pt,outer sep=0pt] (tocancel) {#1};
        \draw[red] (tocancel.south west) -- (tocancel.north east);
    }%
}%
\newcommand\crossout[1]{%
    \setbox0=\hbox{$\displaystyle{#1}$}%
    \hbox{%
	\psline[linecolor=red](0,\ht0)(\wd0,-\dp0)%
	\psline[linecolor=red](0,-\dp0)(\wd0,\ht0)%
	\box0%
	}%
    }
\newcommand\nextline{\unskip\nobreak\hfill\break}
%
\makeatletter
\def\befored@t#1.#2.#3;{#1}
\def\afterd@t#1.#2.#3;{#2}
\def\refhead#1{\edef\next{\ref{#1}}\expandafter\befored@t\next..;}
\def\reftail#1{\edef\next{\ref{#1}}\expandafter\afterd@t\next..;}
\def\lsim{\mathrel{\mathpalette\@versim<}}
\def\gsim{\mathrel{\mathpalette\@versim>}}
\def\@versim#1#2{\vcenter{\offinterlineskip
        \ialign{$\m@th#1\hfil##\hfil$\crcr#2\crcr\sim\crcr } }}
\newcommand\becomes[1]{\mathchoice{\becomes@\scriptstyle{#1}}
   {\becomes@\scriptstyle{#1}} {\becomes@\scriptscriptstyle{#1}}
   {\becomes@\scriptscriptstyle{#1}}}
\def\becomes@#1#2{\mathrel{\setbox0=\hbox{$\m@th #1{\,#2\,}$}%
        \mathop{\hbox to \wd0 {\rightarrowfill}}\limits_{#2}}}
\makeatother
%
%-------------------- Assigning, Commenting & Reviewing -------------
%
\newcommand\assgn[1]{\textcolor{red}{(Assigned to #1)}}
\newcommand\commnt[2]{{\bf Comment by #1} -- \textcolor{brown}{\it #2}}
\newcommand\revwd[2]{{\bf Proposed review by #1} -- \textcolor{blue}{#2}}
%
%--------------Theorems environments -------------------
% \makeatletter
% \newtheoremstyle{nummermitklammern}%
%  {\item[\rlap{\vbox{\hbox{\hskip\labelsep \theorem@headerfont
%   (##2)\ ##1\theorem@separator}\hbox{\strut}}}]}%
%  {\item[\rlap{\vbox{\hbox{\hskip\labelsep \theorem@headerfont
%   (##2)\ ##1 (##3)\theorem@separator}\hbox{\strut}}}]}%
% \makeatother
% \theoremstyle{nummermitklammern}
% \theorembodyfont{\itshape}
% \theoremsymbol{\ensuremath{\square}}
% \newtheorem{formato}{}[section]
% \newtheorem{definition}[formato]{Definition}
% \newtheorem{lemma}[formato]{Lemma}
% \newtheorem{proposition}[formato]{Proposition}
% \newtheorem{theorem}[formato]{Theorem}
% \newtheorem{corollary}[formato]{Corollary}

% \newenvironment{proof}[1][Proof]{\begin{trivlist}
% \item[\hskip \labelsep {\bfseries #1}]}
% {\ensuremath{\blacksquare} \end{trivlist}}
%
\newtheorem{custprop}[theorem]{Proposition}
\newenvironment{remark}[1][Remark:]{\begin{trivlist} \it
\item[\hskip \labelsep {\scshape #1}]}{\end{trivlist}}
%-------------------- formatting-----------------------------------------
%
%\setlength{\textheight}{8.5in} \setlength{\textwidth}{6.5in}
%\setlength{\oddsidemargin}{0in} \setlength{\topmargin}{0in}
%\psset{unit=1mm,linecolor=black,linewidth=1pt}
%\def\normalbaselines{%
%	\normalbaselineskip=20pt plus 0.2pt minus 0.1pt
%	\baselineskip=\normalbaselineskip
%	\lineskip=2pt plus 0.1pt minus 0.1pt
%	\lineskiplimit=2pt
%	}
%\def\fixformat{%
%	\normalbaselines
%	\abovedisplayskip=15pt plus 5pt minus 3pt
%	\belowdisplayskip=\abovedisplayskip
%	\parskip=6pt plus 2pt minus 1pt
%	\skip\footins=20pt plus 10pt minus 5pt
%	\predisplaypenalty=5000
%	\postdisplaypenalty=500
%	\interlinepenalty=50
%	\interdisplaylinepenalty=10000
%	\flushbottom
%	}
%\sectionfont{\huge\rm\blue\nohang\centering}
%\subsectionfont{\Large\sc\nohang\raggedright}
%\makeatletter
%\def\@seccntformat#1{%
%	\@ifundefined{#1@cntformat}%
%	{\csname the#1\endcsname\quad}% default
%	{\csname #1@cntformat\endcsname}% individual control
%	}
%\def\section@cntformat{\thesection.\enspace}
%%\def\subsection@cntformat{\thesubsection:\quad}
%\@addtoreset{equation}{section}
%\newif\iffntmark \fntmarktrue
%\renewcommand\@makefntext[1]{\noindent
%	\iffntmark\llap{\@thefnmark\enspace}\fi
%	#1\unskip }
%\newif\ift@c \t@cfalse
%\def\br@@k{\relax\ift@c\else\unskip\break\fi}
%\def\brk{\protect\br@@k}
%\let\t@c=\tableofcontents
%\renewcommand\tableofcontents{\t@ctrue\t@c\t@cfalse}
%\makeatother
%\renewcommand\theequation{\thesection.\arabic{equation}}
%\renewcommand\thefootnote{{$\langle\arabic{footnote}\rangle$}}

